\documentclass[11pt,a4paper]{article}

\usepackage[margin=1in]{geometry}
\usepackage{amsmath,amssymb,amsthm,mathtools}
\usepackage{booktabs}
\usepackage{enumitem}
\usepackage{microtype}
\usepackage{hyperref}
\usepackage{listings}
\usepackage{xcolor}

\title{\textbf{Recursive Nearby-Factorization and Quotient-Product Descent}}
\author{IamLupo}
\date{30-08-2026}

\newtheorem{proposition}{Proposition}
\newtheorem{definition}{Definition}
\newtheorem{remark}{Remark}

\begin{document}
\maketitle

\begin{abstract}
This paper develops a recursive viewpoint on the factorization of an integer
$n=pq$ using auxiliary factorizations of nearby integers $n+x$. The central
construction is the quotient-product relation induced by two moduli
$r_1,r_2$. Writing
\[
    p=a+k r_1,\qquad q=b+\ell r_2,
\]
leads to the exact quantity
\[
    T=\left\lfloor \frac{n}{r_1r_2}\right\rfloor,
    \qquad
    E=T-k\ell,
\]
and hence
\[
    K=k\ell=T-E.
\]
The important observation is that $K$ is a much smaller integer than
$n$ in the parameter regime of interest. If an auxiliary integer $n+x$
is easier to factor, then its factorization can reveal its own quotient
coordinates $(k_x,\ell_x)$ and therefore its smaller quotient product
$K_x=k_x\ell_x$. This creates a recursive chain in which a large
factorization problem is related to a sequence of smaller arithmetic
problems.

A further simplification is obtained by choosing $x$ so that $n+x$ is even.
For odd semiprime $n$, an odd $x$ gives
\[
    n+x=2m,
\]
so at least one prime factor is known immediately. More generally,
choosing $x$ such that $n+x$ has a prescribed collection of small factors
can produce an auxiliary number whose remaining cofactor is substantially
smaller than the original integer.

The purpose of this work is not to claim a complete factoring algorithm,
but to formalize the recursive mechanism, identify the information carried
by $E_x$, $K_x$, $k_x$, and $\ell_x$, and state the precise reconstruction
problem that a successful recursive method would have to solve.
\end{abstract}

\section{Introduction}

The starting problem is the factorization of a composite integer
\[
    n=pq,
\]
where $p$ and $q$ are unknown factors. Earlier investigations considered
representations based on two moduli and studied whether modular information
could directly identify $p$ and $q$. These investigations showed that
congruences such as
\[
    pq\equiv n\pmod{r_1r_2}
\]
are primarily relational: once one factor is supplied, the other is
immediately determined modulo the chosen modulus. The difficult part is
therefore not verification, but generation of a useful candidate location.

The present approach changes the viewpoint. Instead of attempting to extract
$p$ and $q$ directly from one representation of $n$, we introduce nearby
integers $n+x$ and attempt to factor those auxiliary integers first.
The resulting factors define smaller quotient coordinates and, in turn,
smaller quotient products. This creates a possible recursive descent.

The principal chain is
\[
    n
    \longrightarrow
    n+x
    \longrightarrow
    (p_x,q_x)
    \longrightarrow
    (k_x,\ell_x)
    \longrightarrow
    K_x=k_x\ell_x.
\]
The research question is whether enough information survives this chain to
reconstruct the factor coordinates associated with the original $n$.

\section{Two-Modulus Quotient Coordinates}

Let $r_1,r_2$ be positive integers, typically prime and chosen so that
$\gcd(r_1,r_2)=1$. Write the unknown factors in quotient-remainder form:
\begin{align}
    p &= a+k r_1,\\
    q &= b+\ell r_2,
\end{align}
with
\[
    0\le a<r_1,
    \qquad
    0\le b<r_2.
\]
Substitution into $n=pq$ gives the exact expansion
\begin{equation}
\label{eq:product-expansion}
    n
    =k\ell r_1r_2
     +a\ell r_2
     +bk r_1
     +ab.
\end{equation}

Define
\begin{equation}
\label{eq:T-definition}
    T=\left\lfloor\frac{n}{r_1r_2}\right\rfloor
\end{equation}
and
\begin{equation}
\label{eq:E-definition}
    E=T-k\ell.
\end{equation}
Then the smaller quantity
\begin{equation}
\label{eq:K-definition}
    K=k\ell=T-E
\end{equation}
is an exact integer representation of the quotient coordinates.

The important point is the scale separation. When $r_1r_2$ is large relative
to the relevant factors, $k$ and $\ell$ can be much smaller than $p$ and
$q$, and consequently $K=k\ell$ can be much smaller than $n$.
This does not solve the original problem by itself, because $E$ is unknown.
It does, however, identify a potentially useful intermediate problem:
recover $K$, factor $K$, and use its divisor pairs as candidates for
$(k,\ell)$.

\section{The Defect $E$}

Equation \eqref{eq:product-expansion} gives
\[
    n-k\ell r_1r_2
    =a\ell r_2+bk r_1+ab.
\]
After division by $r_1r_2$, the integer part of the right-hand side creates
the defect between $T$ and $k\ell$.
A useful exact decomposition is obtained by writing
\begin{align}
    a\ell &= c_1r_1+d_1,
        &0\le d_1<r_1,\\
    bk &= c_2r_2+d_2,
        &0\le d_2<r_2.
\end{align}
The resulting carry identity has the form
\begin{equation}
\label{eq:E-carry}
    E=c_1+c_2+c_3,
\end{equation}
where $c_3$ is the final carry generated by the remaining residue terms.
In particular,
\[
    c_1=\left\lfloor\frac{a\ell}{r_1}\right\rfloor,
    \qquad
    c_2=\left\lfloor\frac{bk}{r_2}\right\rfloor.
\]

This decomposition is valuable because it demonstrates that $E$ is not an
arbitrary error term. It is an exact bookkeeping quantity generated by the
cross terms of the product expansion.

Nevertheless, knowledge that $E$ has internal structure is not equivalent
to knowing $E$ from $n$ alone. The recursive strategy described below is
therefore designed to obtain new information from related integers rather
than relying exclusively on a direct inversion of $E$.

\section{Nearby Integers and the $E_x$ Construction}

For an integer displacement $x$, define
\[
    n_x=n+x.
\]
Suppose $n_x$ admits a factorization
\[
    n_x=p_xq_x.
\]
Using the same two moduli, write
\begin{align}
    p_x &= a_x+k_xr_1,\\
    q_x &= b_x+\ell_xr_2.
\end{align}
Define
\begin{align}
    T_x&=\left\lfloor\frac{n_x}{r_1r_2}\right\rfloor,\\
    E_x&=T_x-k_x\ell_x,\\
    K_x&=k_x\ell_x=T_x-E_x.
\end{align}
Thus every nearby factorization provides a complete quadruple
\[
    (T_x,E_x,k_x,\ell_x),
\]
and the smaller product
\[
    K_x=k_x\ell_x.
\]

This is fundamentally different from using $n+x$ merely as another integer
to factor. The nearby factorization is converted into a lower-scale
representation that is directly comparable with the original quotient
coordinates.

\section{Exact Relation Between $E_x$ and $K_x$}

There is an immediate identity relating the original and nearby states:
\begin{equation}
\label{eq:exact-difference}
    E_x-E
    =
    \left(T_x-T\right)-\left(K_x-K\right).
\end{equation}
This equation is exact.

If the quotient $
T$ does not change between $n$ and $n+x$, then
\[
    T_x=T,
\]
and therefore
\begin{equation}
\label{eq:constant-T}
    E_x-E=-(K_x-K).
\end{equation}
Consequently,
\[
    E_x=E
    \quad\Longleftrightarrow\quad
    K_x=K
\]
within any region where $T_x=T$.

This observation is important for the recursive proposal. It shows that
$E_x$ does not need to be treated as an isolated mysterious quantity.
Once the nearby factorization gives $K_x$, its relationship to $T_x$
produces $E_x$ automatically.

\section{Why Choose $n+x$ to Be Even?}

Assume the original integer $n$ is odd, as occurs when $p$ and $q$ are odd
primes. If $x$ is odd, then
\[
    n+x
\]
is even. Hence
\begin{equation}
\label{eq:even-neighbor}
    n+x=2m,
    \qquad
    m=\frac{n+x}{2}.
\end{equation}
The prime factor $2$ is therefore known without any search.

More generally, if $x$ is chosen so that
\[
    n+x=2^t u,
\]
then all factors of $2^t$ are known immediately and only the odd cofactor
$u$ remains to be factored.

The recursive viewpoint is not restricted to powers of two. One can seek
auxiliary displacements for which
\[
    n+x=s\,u
\]
with $s$ having a known factorization composed of small primes. Then the
remaining factorization problem concerns the smaller cofactor $u$.

The crucial distinction is that the auxiliary integer is not expected to
be useful merely because it is numerically close to $n$. It is useful when
its factorization is deliberately easier or contains known structure.

\section{The Recursive Descent Hypothesis}

The proposed recursive process can be expressed as follows.

\subsection{Level 0: Original Integer}

Start with
\[
    n=pq.
\]
The factors $p,q$ are unknown.

\subsection{Level 1: Nearby Auxiliary Integer}

Choose $x$ so that $n+x$ is easier to factor. Compute
\[
    n+x=p_xq_x.
\]
The auxiliary factors define
\[
    k_x=\left\lfloor\frac{p_x}{r_1}\right\rfloor,
    \qquad
    \ell_x=\left\lfloor\frac{q_x}{r_2}\right\rfloor.
\]
Hence
\[
    K_x=k_x\ell_x.
\]

\subsection{Level 2: Smaller Quotient Product}

The number $K_x$ is much smaller than the original $n$ in the intended
regime. Therefore we factor $K_x$ directly:
\[
    K_x=d_1d_2.
\]
Every divisor pair produces a candidate quotient pair
\[
    (k_x,\ell_x)=(d_1,d_2).
\]
The quotient coordinates then define intervals for the corresponding
factor residues:
\begin{align}
    d_1r_1&\le p_x<(d_1+1)r_1,\\
    d_2r_2&\le q_x<(d_2+1)r_2.
\end{align}

\subsection{Level 3 and Beyond}

The same process may be applied recursively to the smaller objects that
arise from the quotient coordinates or from related nearby factorizations.
Schematically,
\begin{equation}
\label{eq:recursive-chain}
    n
    \rightarrow
    n+x_1
    \rightarrow
    K_1
    \rightarrow
    n+x_2
    \rightarrow
    K_2
    \rightarrow\cdots.
\end{equation}
The hypothesis is that repeated descent eventually reaches integers small
enough that their complete factorization is inexpensive. Their recovered
quotient structure is then propagated upward.

\section{What Information Must Propagate Upward?}

The recursive proposal is only useful if information from a lower level
can constrain a higher level. Several quantities are available for this
purpose:

\begin{enumerate}[label=(\arabic*)]
    \item The exact lower-level quotient product $K_x$.
    \item The divisor pairs of $K_x$.
    \item The quotient coordinates $(k_x,\ell_x)$.
    \item The defect $E_x=T_x-K_x$.
    \item The carry decomposition of $E_x$.
    \item The quotient-cell locations implied by $(k_x,\ell_x)$.
    \item Relations between nearby states when $T_x$ remains constant.
\end{enumerate}

The essential reconstruction question is therefore not
\emph{whether} $K_x$ can be factored. It is whether the resulting divisor
pairs provide a sufficiently small set of possibilities for the original
$(k,\ell)$.

\section{The Role of Stability}

Nearby factorizations can sometimes remain in the same quotient cells. In
that case
\[
    k_x=k,
    \qquad
    \ell_x=\ell,
\]
which immediately gives
\[
    K_x=K.
\]
This is a geometric stability phenomenon.

A stronger and more interesting possibility is a cross-cell collision:
\[
    (k_x,\ell_x)\ne(k,\ell)
\]
but
\[
    k_x\ell_x=k\ell.
\]
Such a relation would show that the smaller product $K$ can remain stable
even while the individual quotient coordinates move.

These two cases must be distinguished carefully. Same-cell stability is
expected from quotient geometry; cross-cell equality would provide
additional arithmetic structure.

\section{Candidate Reconstruction of the Original Factors}

Suppose a recursive branch has produced a candidate original quotient pair
$(k,\ell)$. Then the unknown residues satisfy
\[
    p=kr_1+a,
    \qquad
    q=\ell r_2+b,
\]
with
\[
    0\le a<r_1,
    \qquad
    0\le b<r_2.
\]
The exact product equation becomes
\begin{equation}
\label{eq:reconstruction}
    n
    =(kr_1+a)(\ell r_2+b).
\end{equation}
The recursive information reduces the problem to finding a compatible
residue pair $(a,b)$ inside a bounded rectangle.

The carry identity provides additional constraints. In particular,
\[
    c_1=\left\lfloor\frac{a\ell}{r_1}\right\rfloor,
    \qquad
    c_2=\left\lfloor\frac{bk}{r_2}\right\rfloor,
\]
and
\[
    E=c_1+c_2+c_3.
\]
Thus a recovered value of $K$ narrows $(k,\ell)$, while a recovered value
of $E$ narrows the admissible carry states.

The intended chain is therefore
\[
    K
    \rightarrow
    (k,\ell)
    \rightarrow
    (c_1,c_2,c_3)
    \rightarrow
    (a,b)
    \rightarrow
    (p,q).
\]

\section{Why the Recursive Viewpoint Is Different}

The earlier modular experiments primarily used congruences as filters for
already-generated candidates. In contrast, the recursive hypothesis tries
to change the source of the candidates themselves.

The distinction can be summarized as
\begin{align*}
    \text{filtering approach:}
    &\qquad p\text{-candidates}\rightarrow\text{modular tests},\\
    \text{recursive approach:}
    &\qquad n+x\rightarrow\text{smaller factorization}\rightarrow K_x
    \rightarrow\text{quotient candidates}.
\end{align*}

The latter can only be considered a genuine reduction if the auxiliary
factorizations are collectively cheaper than factoring $n$ directly and if
the information they provide can be propagated upward without recreating
the original search space.

\section{The Central Bottleneck}

There are two distinct unknowns in the recursive construction.

First, one must choose useful displacements $x$. An arbitrary nearby number
is not necessarily easier to factor. The auxiliary search therefore needs
a principled way of finding displacements for which $n+x$ has favorable
factor structure.

Second, even when $n+x$ has been successfully factored, the resulting
$K_x$ must identify useful information about the original factorization.
Knowing a small integer is not automatically equivalent to knowing which
of its divisor pairs corresponds to the original quotient coordinates.

Thus the recursive method must eventually establish a statement of the form
\[
    \text{easy auxiliary factorization}
    \quad\Longrightarrow\quad
    \text{small candidate set for }(k,\ell)
\]
and then
\[
    \text{small candidate set for }(k,\ell)
    \quad\Longrightarrow\quad
    \text{small candidate set for }(p,q).
\]

\section{A Precise Recursive Research Program}

A complete experimental program based on this hypothesis should treat the
following as separate stages.

\subsection{Auxiliary Selection}
Find $x$ for which $n+x$ has a known or deliberately favorable factor
structure, beginning with the guaranteed factor $2$ obtained from odd
$x$ when $n$ is odd.

\subsection{Auxiliary Factorization}
Factor the remaining cofactor of $n+x$ and record the exact tuple
\[
    (T_x,E_x,k_x,\ell_x,K_x).
\]

\subsection{Quotient-Product Factorization}
Factor $K_x$ and enumerate its divisor pairs. Since $K_x$ is much smaller
than $n$, this is the intended lower-level factoring step.

\subsection{Upward Reconstruction}
Use the candidate quotient coordinates, carry constraints, and residue
relations to generate candidate factors for the original $n$.

\subsection{Recursive Propagation}
Apply the same mechanism to smaller auxiliary quantities whenever their
factorization exposes another quotient-product representation.

\section{Necessary Conditions for a Successful Method}

A successful recursive mechanism would need to satisfy all of the following
conditions simultaneously:

\begin{enumerate}[label=(\arabic*)]
    \item The chosen $n+x$ values must be significantly easier to factor
    than $n$ itself.
    \item Their factorizations must reveal exact lower-level quotient data.
    \item Factoring the corresponding $K_x$ values must produce a small
    collection of plausible quotient pairs.
    \item The ambiguity among divisor pairs must shrink rather than remain
    proportional to the original search space.
    \item The reconstruction of $(p,q)$ must not require enumerating the
    original factor interval.
    \item The recursive depth needed to reach easily factorable values must
    remain computationally manageable.
\end{enumerate}

Failure of any one of these conditions can turn the apparent recursion into
another representation of the same search problem.

\section{Discussion}

The new viewpoint does not rely on the assumption that a single modular
congruence directly exposes the hidden factors. Instead, it treats
factorization as a hierarchy of related problems.

The basic hierarchy is
\[
    n
    \longrightarrow
    n+x
    \longrightarrow
    (k_x,\ell_x)
    \longrightarrow
    K_x
    \longrightarrow
    \operatorname{factor}(K_x).
\]
The hope is that each step reduces the effective size of the unknown
object and that the resulting information can be propagated back upward.

The special choice of even neighbors is particularly attractive because it
provides an immediate nontrivial factor. For odd $n$, every odd displacement
$x$ gives an even auxiliary integer. This creates a systematic family of
related factorizations rather than requiring an arbitrary second composite.

However, the decisive mathematical question remains whether the recursive
information is genuinely new. In particular, if obtaining and processing
the lower-level states implicitly requires a search whose cost scales like
the original factorization problem, then the method has changed the
representation without changing the underlying difficulty.

The purpose of the recursive framework is therefore to locate a possible
information bottleneck that is different from direct candidate filtering:
can deliberately easy factorizations of related integers reveal the hidden
coordinates of the original factorization?

\section{Conclusion}

The two-modulus quotient representation introduces a small intermediate
quantity
\[
    K=k\ell=T-E,
\]
which is potentially far smaller than the original composite $n$.
The defect $E$ has an exact carry decomposition, but direct inversion of
that decomposition does not automatically identify $K$.

The recursive nearby-factorization hypothesis proposes another route. By
factoring suitably chosen integers $n+x$, one can obtain exact values
$T_x$, $E_x$, $k_x$, and $\ell_x$, and hence the smaller product
\[
    K_x=k_x\ell_x=T_x-E_x.
\]
The factorization of $K_x$ can then generate candidate quotient coordinates.
If suitable auxiliary integers can be found recursively, the factorization
of the original $n$ may be decomposed into a sequence of smaller
factorization problems.

Choosing $n+x$ to be even gives an especially simple first step:
\[
    n+x=2m,
\]
with the factor $2$ known immediately. More generally, auxiliary values
with intentionally favorable small-factor structure may create deeper
recursive reductions.

At present, this framework should be regarded as a precise research
hypothesis rather than a proof of a new factoring algorithm. The essential
future test is whether the information obtained from easier auxiliary
factorizations produces a genuinely shrinking candidate set for the
original $(p,q)$ without recreating the original search.

\end{document}
